% META: {"id": "TNSI-STRUCT-COURS-11", "chapitre": "TNSI-STRUCTURES-DONNEES", "type_objet": "cours", "capacites_codes": ["C11"], "sources_inspiration": [], "mode_creation": "ex_nihilo", "status": "approved", "capacites": ["TNSI-STRUCTURES-DONNEES-C11"]}
\section{Interface et implémentation d'une structure de données}

% ============================================================
% STRATE 1 — Definitions et exemple contraste
% ============================================================

\definition[1]{
L'\textbf{interface} d'une structure de donnees est l'ensemble des \textbf{operations} qu'elle propose (ce que l'on peut faire), \textbf{sans preciser comment} elles sont realisees. Une \textbf{implementation} est une facon concrete de coder ces operations (avec quelles donnees internes, quel algorithme).
}

\margeAppui{
Une meme interface peut avoir \textbf{plusieurs implementations} differentes : l'utilisateur de la structure n'a pas besoin de connaitre l'implementation pour utiliser l'interface.
}

\exemple{
\textbf{Interface d'une pile (Stack)} : \lstinline{empiler(valeur)}, \lstinline{depiler()}, \lstinline{est_vide()}.
}

\textbf{Implementation 1 — avec une liste Python (fin de liste = sommet) :}
\begin{python}
class PileListe:
    def __init__(self):
        self._donnees = []

    def empiler(self, valeur):
        self._donnees.append(valeur)

    def depiler(self):
        return self._donnees.pop()

    def est_vide(self):
        return len(self._donnees) == 0
\end{python}

\textbf{Implementation 2 — avec une liste chainee (debut = sommet) :}
\begin{python}
class Maillon:
    def __init__(self, valeur, suivant=None):
        self.valeur = valeur
        self.suivant = suivant

class PileChainee:
    def __init__(self):
        self._sommet = None

    def empiler(self, valeur):
        self._sommet = Maillon(valeur, self._sommet)

    def depiler(self):
        v = self._sommet.valeur
        self._sommet = self._sommet.suivant
        return v

    def est_vide(self):
        return self._sommet is None
\end{python}

% BEGIN-VERIFY
% class PileListe:
%     def __init__(self):
%         self._donnees = []
%     def empiler(self, valeur):
%         self._donnees.append(valeur)
%     def depiler(self):
%         return self._donnees.pop()
%     def est_vide(self):
%         return len(self._donnees) == 0
%
% class Maillon:
%     def __init__(self, valeur, suivant=None):
%         self.valeur = valeur
%         self.suivant = suivant
%
% class PileChainee:
%     def __init__(self):
%         self._sommet = None
%     def empiler(self, valeur):
%         self._sommet = Maillon(valeur, self._sommet)
%     def depiler(self):
%         v = self._sommet.valeur
%         self._sommet = self._sommet.suivant
%         return v
%     def est_vide(self):
%         return self._sommet is None
%
% # meme interface, deux implementations : meme comportement observable
% for Pile in (PileListe, PileChainee):
%     p = Pile()
%     assert p.est_vide()
%     p.empiler(1)
%     p.empiler(2)
%     p.empiler(3)
%     assert not p.est_vide()
%     assert p.depiler() == 3
%     assert p.depiler() == 2
%     assert p.depiler() == 1
%     assert p.est_vide()
% END-VERIFY

\exemple{
Ces deux implementations offrent \textbf{exactement la meme interface} (memes trois methodes, meme comportement observable), mais des structures de donnees internes totalement differentes (tableau dynamique contre liste chainee).
}

% ============================================================
% STRATE 2 — Erreurs frequentes
% ============================================================

\erreurFrequente{
\textbf{Confondre interface et implementation lors d'un test.} Un test doit porter sur le comportement observable via l'interface (par exemple : apres avoir empile 1, 2, 3, depiler doit renvoyer 3), pas sur les details internes (comme la structure exacte de \lstinline{_donnees}). Cela permet au meme jeu de tests de valider plusieurs implementations.
}
